\section{Divergence as Boundary Signal}

\label{sec:divergence}

From the preceding discussion,
one conclusion becomes unavoidable.
Physics operates on registered differences.
Its equations describe quantities that have become stable enough to be
compared,
measured,
or related inside a generated domain.

A divergence should therefore not be interpreted too quickly as the
literal appearance of an actual infinite object.
It may indicate that the variables of a theory are being asked to
represent differences beyond the scale,
texture,
or reference structure in which those variables remain finite.

\medskip
\begin{quote}
Divergence is not simply an infinite object appearing inside theory.
\end{quote}
\medskip

Nor is divergence merely a failure of rigor.
In many cases,
a divergence signals that a descriptive window has been pushed beyond
the differences it can stably register.
The theory remains meaningful within its domain,
but its variables may no longer be able to represent the required
differences as finite quantities.

\subsection*{Registered Difference and Overextension}

The point is structural.
Physical description depends on traces,
scales,
and quantities that have already become describable.
When a theory compresses difference toward a point,
amplifies difference toward unbounded energy,
or iterates difference beyond the stability of its own variables,
divergence can appear.

From the perspective of \VED{},
divergence is read as a boundary signal:
the trace of a description being pushed beyond the differences it can
stably register.

\medskip
\begin{quote}
Divergence marks descriptive overextension.
\end{quote}
\medskip

This statement should not be overread.
\VED{} does not claim that all divergences have the same physical
origin.
Different theories,
scales,
and mathematical structures produce divergences in different ways.
The claim is interpretive and structural:
divergences can indicate that a successful descriptive window has
approached a local boundary of finite representation.

\subsection*{What Divergence Does Not Reveal}

Divergence does not turn the boundary into an object of description.
It does not provide access to a hidden region,
nor does it turn the \ungenerated{} into something observable.
It reflects the limit of the current description.

\medskip
\begin{quote}
Divergence does not cross the boundary.\\
It marks the limit of a window.
\end{quote}
\medskip

This is why divergence cannot be treated only as a value to be
discarded.
It may also be a diagnostic sign.
It indicates that the theory has reached the edge of the differences
its own variables can carry as finite quantities.

\subsection*{Implication for Renormalization}

At this point,
the meaning of renormalization begins to shift.
If divergence marks boundary contact or descriptive overextension,
then renormalization is not simply the deletion of a mistake.
It is the finite re-expression of that contact inside a chosen
observational or mathematical window.

\medskip
\begin{quote}
Divergence signals boundary contact;\\
renormalization re-expresses that contact finitely.
\end{quote}
\medskip

This will be examined in detail in the next section.
